% ===== Opening: title, hook, the book and the series, today's plan (~5 min) =====

\webinartitle{2}{Գլուխ 2. Հիշողության նվերը}%
  {\chip{door}{կարճ ու երկար} \chip{eyes}{դեմք ու անուն} \chip{wine}{ալկոհոլ} \chip{crown}{եսասեր հիշողություն} \chip{magnifier}{երբ այն փչանում է}}%
  {thinking}{Ֆայլը գտնվեց։ Մի քիչ խմբագրված։}{{door}{eyes}{crown}}

% Hook: four everyday memory glitches; the audience types the ones they know
\begin{frame}{Ծանո՞թ է}
  \begin{columns}[c]
    \begin{column}{0.58\textwidth}
      \begin{popup}{ուղեղ.exe}
        {\blackfont\large Ֆայլը չի գտնվել}\par\medskip
        \begin{enumerate}
          \item<2-> Մտար սենյակ ու մոռացար՝ ինչու
          \item<3-> Դեմքը ճանաչում ես, անունը՝ ոչ
          \item<4-> Բառը լեզվիդ ծայրին է, բայց չի գալիս
          \item<5-> Խնջույքից հետո չես հիշում՝ ինչպես հասար տուն
        \end{enumerate}
      \end{popup}
    \end{column}
    \begin{column}{0.36\textwidth}
      \centering
      \mascot[3.2cm]{confused}\par\medskip
      \onslide<6->{%
      \begin{chatbox}
        Գրիր համարները, որոնք քեզ հետ եղել են (օր.՝ \textbf{1, 3})
      \end{chatbox}}
    \end{column}
  \end{columns}
\end{frame}

% The book and where this chapter sits in the series (for newcomers)
\begin{frame}{Գիրքը, հեղինակը և մեր շարքը}
  \begin{columns}[T]
    \begin{column}{0.5\textwidth}
      \begin{card}
        {\blackfont\Large The Idiot Brain} {\small (2016)}\par
        «Ապուշ ուղեղը»\par\medskip
        \begin{itemize}
          \item \textbf{Դին Բըրնեթ}՝ նեյրոգիտնական, Քարդիֆի համալսարան
          \item Guardian-ի «Brain Flapping» բլոգը, ստենդափ
          \item Գրված է կատակով, բայց գիտական աղբյուրներով
        \end{itemize}
        \vspace{-1mm}
      \end{card}
    \end{column}
    \begin{column}{0.44\textwidth}
      {\blackfont\small Գրքի 8 գլուխը՝ 8 վեբինարում}\par\smallskip
      {\color{plum!55!cream}\small \emoji{check}\ Ներածություն, 1՝ վեբինար 1}\par\smallskip
      \chip[lavender]{thought}{2. Հիշողություն. այսօր}\par\smallskip
      {\color{plum!55!cream}\small
      3. Վախ \quad 4. Խելք \quad 5. Ընկալում\\
      6. Անհատականություն \quad 7. Ուրիշ մարդիկ\\
      8. Երբ ուղեղը «փչանում» է}\par\smallskip
      \begin{card}[colback=mint!35!cardwhite]
        Վեբինար 1-ին չե՞ս եղել։ Ոչինչ. ամեն մեկն ինքնուրույն է։
      \end{card}
    \end{column}
  \end{columns}
\end{frame}

% Plan for today: the five topics of chapter 2, split by the break
\begin{frame}{Այսօր}
  \begin{columns}[T]
    \begin{column}{0.4\textwidth}
      \begin{card}[colback=lavender!45!cream]
        {\blackfont Մինչև ընդմիջումը}\par\smallskip
        \chip{door}{կարճ ու երկար}\par\smallskip
        \chip{eyes}{դեմքեր ու անուններ}
      \end{card}
    \end{column}
    \begin{column}{0.4\textwidth}
      \begin{card}[colback=mint!45!cream]
        {\blackfont Ընդմիջումից հետո}\par\smallskip
        \chip{wine}{ալկոհոլ}\par\smallskip
        \chip{crown}{եսասեր հիշողություն}\par\smallskip
        \chip{magnifier}{երբ այն փչանում է}
      \end{card}
    \end{column}
  \end{columns}
  \vspace{1mm}
  \centering
  {\small \emoji{coffee}\ Մեջտեղում՝ 5 րոպե ընդմիջում \quad \emoji{speech}\ Հարցերը գրեք չաթում ցանկացած պահի}
\end{frame}
